\chapter{Modeling Constant Function Market Makers (CFMM)}
A new form of financial market, Decentralized Exchanges (DEXs) also referred
to as pool, have found its roots on top of blockchain technology. As in any
financial market, DEXs allow an agent to trade assets, in this case
cryptocurrencies (digital assets). Thorough a decentralized ledger stored on
multiple computers (blockchain) the DEXs keeping track of ownership by
storing in their storage slot. This report goes through the underlying
principles of DEXs and their properties by mathematical modeling and
optimization.

\section{Basics}
A DEX is governed by a constant function market maker (CFMM), which is a
function $\varphi:\mathbb{R}^{n}_{+} \to \mathbb{R}_+$ the trading function
and a set of conditions. Here $n \ge 2$ is the number of assets an agent can
trade, these are labeled $\{1,\ldots, n\}$. The CFMM is a function of the
reserves $y \in \mathbb{R}^{n}_+$, where $y_i \in \mathbb{R}_+$ is the
reserve of asset $i \in \{1,\ldots,n\}$, these are in the portfolio of the
DEX. There are two basic ways an agent can interact with the DEX:
\begin{itemize}
    \item Change liquidity
    \item Trade
\end{itemize}
This report covers the agents' iteration by trade. The liquidity change
could be added in the final master thesis, as part of a general introduction
the CFMMs.

An agent initializes a trade on a DEX, in which a set amount of assets is
traded for another. A set amount of assets given is denoted by $x \in
\mathbb{R}^{n}_+$, and the set amount of assets received by $\tilde{x} \in
\mathbb{R}^{n}_+$. Then $x_i, \tilde{x}_i \in \mathbb{R}_+$ is the amount of
given/received from asset $i \in \{1,\ldots,n\}$ respectively. A trade can either be
accepted or rejected, depending on an acceptance principle of the underlying
CFMM. If the trade is accepted the state of the pool is updated by changing
the reserves
\begin{align}
    y \mapsto y + x - \tilde{x}.
\end{align}
A necessary condition for trade acceptance is $y \ge 0$ after the trade. In
the case the trade is rejected the reserves do not change. Throughout
assume that the vectors $x$ and $\tilde{x}$ have disjoint support. This
is a valid assumption, since when trading, the agent is never giving and receiving
the same asset, but rather one or multiple different ones. The
support of a vector $v \in \mathbb{R}^{n}_+$ is $\text{supp}(v) := \{i\; |\;
v_i > 0\}$, thus the assumption is $\text{supp}(x)\;\cap\;
\text{supp}(\tilde{x}) = \emptyset$.

A trade $(x, \tilde{x})$ is accepted, if the trading function $\varphi$
stays constant during the trade
\begin{align}
    \varphi(y+\gamma x - \tilde{x}) = \varphi(y).
\end{align}
Depending on the CFMM, a trade has a fee $\gamma \in (0, 1]$, where
$(1-\gamma)x$ of the set amount given assets is distributed to liquidity
providers for the reward of providing liquidity. The discounted amount
$\gamma x$ is actually being exchanged for $\tilde{x}$, but the reserves are
updated by the full amount $x - \tilde{x}$.

Throughout assume that the trading function $\varphi$ is concave,
increasing and differentiable, as it is implemented by Uniswap
\cite{uniswap}, Balancer \cite{balancer}, Sushiswap \cite{sushiswap}, CurveFI
\cite{curvefi} and many others. Most of them are also positive homogenious
$\varphi(\alpha y)=\alpha(y)$, for all $\alpha \ge 0$.

\subsection{CFMMs Examples}
The weighted sum trading function is
\begin{align}
    \varphi(y) = p^{T}y = \sum_{i=1}^{n} p_i y_i, \label{eq: weighted_sum}
\end{align}
where $p_i$ is called the price of asset $i$. The trade acceptance condition
simplifies to
\begin{align}
    &\varphi(y+\gamma x- \tilde{x}) = \varphi(y) \nonumber \\
    &p^{T} (y +\gamma x - \tilde{x}) = p^{T} y \nonumber \\
    &\gamma p^{T} x = p^{T} \tilde{x}.
\end{align}
The weighted geometric mean trading function
\begin{align}
    \varphi(y) = \prod_{i=1}^{n} y_i^{w_i}, \label{eq: weighted_geo_mean}
\end{align}
where $w_i > 0 $ and $\sum_{i=1}^{n} w_i=1$. The trade acceptance condition can be
simplified by splitting the product ($\ln$ of sum) in the indices of
$\text{supp}(x)$ and $\text{supp}(\tilde{x})$
\begin{align}
    \varphi(y+\gamma x- \tilde{x}) &= \varphi(y) \nonumber \\
    \prod\limits_{\substack{i \in \text{supp}(x) \\
    j \in \text{supp}(\tilde{x}) \\
    k \notin \text{supp}(x)\; \cup\; \text{supp}(\tilde{x})  }}
    (y_i + \gamma x_i)^{w_i}
    (y_j - \tilde{x}_j)^{w_j}
    (y_k)^{w_k} &= \prod_{i=1}^{n} (y_i)^{w_i} \nonumber\\
    \prod\limits_{\substack{i \in \text{supp}(x) \\
    j \in \text{supp}(\tilde{x})}}
    \left(1 + \frac{\gamma x_i}{y_i}\right)^{w_i}
    \left(1 - \frac{\tilde{x}_j}{y_j}\right)^{w_j}
    &= 1.\label{eq: wm_trade_acceptance}
\end{align}
For $|\text{supp}(x)| = |\text{supp}(\tilde{x})| = 1$, let $i$ and $j$ be the
corresponding indices of the giving and receiving asset then an analytic
expression can be explicitly derived  for $x_i$ (or $\tilde{x}_j$) by
taking the power of $w_i$ (or $w_j$) and rearranging of equation
\ref{eq: wm_trade_acceptance}.
\begin{align}
&x_i = \frac{y_i}{\gamma}\left(
\frac{y_j^{w_j/w_i}}{\left(y_j - \tilde{x}_j \right)^{w_j/w_i}} - 1
\right) \label{eq: gm_fef}\\
&\tilde{x}_j= y_j\left(
    1 - \frac{y_i^{w_i/w_j}}{\left(y_i+\gamma x_i\right)^{w_i/w_j} }
\right) \label{eq: gm_ref}
\end{align}
Note that the weighted geometric mean is positive homogenious because the
weights sum up to one, i.e. $\sum_{i=1}^{n}w_i = n $


Curve \cite{stableswap} market makers implement the following trading function
\begin{align}
    \varphi(y) =\sum_{i=1}^{n}y_i  - \alpha \prod_{i=1}^{n} y_i^{-1} ,
\end{align}
for $\alpha > 0$, used for stable coin markets to reduce price impact without
large counter liquidity. An aggregated generalization is a combination between
the geometric weighted mean and the (un)weighted sum
\begin{align}
    \varphi(y) =(1-\alpha)\sum_{i=1}^{n}y_i  - \alpha \prod_{i=1}^{n} y_i^{w_i},
\end{align}
where $\| w\|_1 = 1$, $w_i > 0$ and $\alpha \in [0,1]$. For $\alpha = 0$, the
trading function becomes the weighted sum and for $\alpha = 1$ the geometric
weighted mean.

\subsection{Concentrated Liquidity}
Another implementation of CFMM are automated market makers with concentrated
liquidity like UniswapV3 \cite{uniswapv3_paper}. The idea is to split
liquidity provision from the whole sublevel set
$\varphi(y +\gamma x-\tilde{x}) =\varphi(y)$ to multiple sets. The splitting
happens is determined by ticks $\{p(i)\}_{i \in \mathbb{Z} }$ and the rate
usually set $p(i) = 1.0001^{i}$, where $p(i)$ is the so-called marginal
exchange rate of the tick. Every tick range
\begin{align}
    (p(i), p(i+1)],
\end{align}
has its own part of reserves (liquidity) and is governed by a one trading
function $\varphi$ depending on the liquidity bound by the tick range. The
trade is executed, when the marginal exchange rate is in the rage, on a
specifically provided liquidity. The liquidity providers specify an upper and
a lower bound $(p(U), p(L)]$ and provide liquidity on depth $\tilde{k} =
\varphi(y)$. They earn, split by weighted amount liquidity provided, the fees
of the trade associated with the CFMM.

During a trade when the marginal exchange rate raises above a trick range,
the tick range is flipped to the next one and the rest of the trade is
executed on the new tick range. This however is associated with more
transaction cost for the agent conducting the trade.
\section{First order approximations\label{sec: first_order_approx}}
Let $\nabla \varphi(y) = P$ be the vector of unscaled prices, $P$ is positive
in all entries because $\varphi$ is by assumption increasing. The entries by
the gradient are
\begin{align}
    P_i = \frac{\partial \varphi(y)}{\partial y_i}
\end{align}
The reported price also called the marginal exchange rate is defined as
ration between the unscaled prices and the last unscaled price
\begin{align}
    p_i = \frac{P_i}{P_n},
\end{align}
where $p_n = 1$ always.

Recall the trade acceptance condition
\begin{align}
    \varphi(y + \gamma x - \tilde{x}) - \varphi(y) = 0.
\end{align}
First order Taylor expansion of $(y + \gamma x - \tilde{x})$ around $y$ gives us
\begin{align}
    \varphi(y + \gamma x - \tilde{x}) - \varphi(y) &\simeq \nabla
    \varphi(y)^{T} (\gamma x - \tilde{x}) \\
                   &=P^{T} (\gamma x - \tilde{x}),
\end{align}
rearranging gives
\begin{align}
    \gamma \sum_{i \in \text{supp}(x) } P_i x_i \simeq \sum_{i \in
    \text{supp}(\tilde{x}) } P_i \tilde{x}_i,
\end{align}
which is positive homogenious. Hence, dividing by $P_n$
\begin{align}
    \gamma \sum_{i \in \text{supp}(x) } p_i x_i \simeq \sum_{i \in
    \text{supp}(\tilde{x}) } p_i \tilde{x}_i,
\end{align}
also holds.
Note that this is a first order approximation, similarly to a
midpoint price in an order book. The actual trades are completely
determined by the trade acceptance condition.
\section{Swaps}
The most common type of trade interaction with a DEX is a swap, i.e. a two
asset trade. Denote these two assets with index $i$ and index $j$, then
$x=\delta e_i$, $\tilde{x}=\tilde{\delta} e_j$, where $\delta, \tilde{\delta}
\in \mathbb{R}_+$ and $e_i,e_j$ are the canonical basis vectors of
$\mathbb{R}^{n}$, $i \neq j$ (disjoint support). The trade acceptance
condition is then
\begin{align}
    \varphi(y + \gamma\delta e_i - \tilde{\delta}e_j) = \varphi(y).
    \label{eq: trade_acceptance_condition}
\end{align}
It follows from the assumptions on $\varphi$, that the left-hand side of the
equation $\varphi(y + \gamma\delta e_i - \tilde{\delta}e_j)$ is increasing
in $\delta$ and decreasing in $\tilde{\delta}$. Hence, for all $\delta \ge 0$
there is at most one $\tilde{\delta}$ which satisfies the trade acceptance
condition. Respectively for all $\tilde{\delta} \ge 0$ there is at most one
$\delta$ satisfying the trade acceptance condition. Thereby $\delta$ and
$\tilde{\delta}$ have a one-to-one function correspondence, and two asset
trades are uniquely determined by either $\delta$ or $\tilde{\delta}$. By
concavity we get that
\begin{align}
    0 &= \varphi(y+\gamma\delta e_i - \tilde{\delta}e_j) - \varphi(y)\\
      &\le \nabla_y \varphi(y) (\gamma\delta e_i - \tilde{\delta}e_j) \\
      &\le p_i \gamma \delta - p_j \tilde{\delta}.
\end{align}
Thus that
\begin{align}
    p_j \tilde{\delta} \le p_i \gamma \delta.
\end{align}

In this case it is natural to define this function $F: \mathbb{R}_+ \to
\mathbb{R}$, and call it the forward exchange function. Where $F(\delta)$ is
unique to one $\tilde{\delta}$. The function $F$ is objective for all valid
trades, additionally the function is increasing since $\varphi$ is
increasing. Also, $F$ is non-negative, because $F(0) = 0$ and it is concave.
Use the implicit function theorem on the trade acceptance condition
\begin{align}
f(\delta, \tilde{\delta}) = \varphi(y + \gamma\delta e_i -\tilde{\delta}e_j)
- \varphi(y) =0,
\end{align}
denote $y' = y+\gamma\delta e_i + \tilde{\delta}e_j$, the chain rule gives
\begin{align}
    F'(\delta) &= -\frac{\partial_\delta f(\delta,
    \tilde{\delta})}{\partial_{\tilde{\delta}} f(\delta, \tilde{\delta})}\\
       &=\gamma \frac{\partial_{y_i}\varphi(y')}{\partial_{y_j}
       \varphi(y')} \label{eq: forward_implicit}
\end{align}
Together with the properties of $\varphi$, it can be shown that $F$ is
concave by showing that
\begin{align}
    F(\delta') \le F'(\delta)(\delta' - \delta) + F(\delta),
\end{align}
holds for all $\delta, \delta' \ge 0$ valid. Let $y'' = y + \gamma \delta' e_i -
F(\delta') e_j$ and note that $\varphi(y) = \varphi(y') = \varphi(y'')$ as
part of the trade acceptance condition. Since $\varphi$ is a concave function
\begin{align}
    \varphi(y'') \le \nabla_y\varphi(y')^{T} (y'' -y')+
    \varphi(y'),\label{eq: phi_concave}
\end{align}
or simply
\begin{align}
    \nabla_y\varphi(y')^{T} (y'' - y') \ge 0.
\end{align}
Substituting for $y''$ and $y'$ gets
\begin{align}
    &0 \le (\delta' -\delta)\partial_{y_i} \varphi(y') - (F(\delta') -
    F(\delta))\partial_{y_j}\varphi(y') \\
    &0 \le (\delta' -\delta)\frac{\partial_{y_i} \varphi(y')}{\partial_{y_j}
    \varphi(y')} - (F(\delta') - F(\delta))\\
    & F(\delta') \le F'(\delta)(\delta' - \delta) + F(\delta').
\end{align}

Respectively a definition of the reverse exchange function
$G: \mathbb{R}_+ \to \mathbb{R} \cup {\infty}$, $G(\tilde{\delta})$ uniquely determined by
$\delta$ satisfying the trade acceptance condition. If no such $\delta$
exists then the function is mapped to $\{\infty\}$. The reverse exchange
function is non-negative and increasing because of $\varphi$, but it is
convex. To show convexity, use the implicit function theorem
\begin{align}
    G'(\tilde{\delta})  &= -\frac{\partial_{\tilde{\delta}} f(\delta,
    \tilde{\delta})}{\partial_{\delta} f(\delta, \tilde{\delta})}\\
                        &=\frac{1}{\gamma} \frac{\partial_{y_j}\varphi(y')}{
                    \partial_{y_i} \varphi(y')}.
\end{align}
It needs to be shown that
\begin{align}
    G(\tilde{\delta}') \ge G'(\tilde{\delta})(\tilde{\delta}' -
    \tilde{\delta}) + G(\tilde{\delta}),
\end{align}
for all $\tilde{\delta},\tilde{\delta}' \ge 0$. Set $y'' = y + \gamma
G(\tilde{\delta}')e_i - \tilde{\delta}'e_j$ similarly like with the forward
exchange function use the concavity of the trading function $\varphi$ and
the trade acceptance condition
\begin{align}
    &\nabla_y \varphi(y')^{T} (y'' - y') \ge 0\\
    &(G(\tilde{\delta}') - G(\tilde{\delta}))\partial_{y_i}\varphi(y')
    - (\tilde{\delta}' - \tilde{\delta})\partial_{y_j}\varphi(y') \ge 0\\
    &(G(\tilde{\delta}') - G(\tilde{\delta}))
    - (\tilde{\delta}' - \tilde{\delta})\frac{\partial_{y_j}
    \varphi(y')}{\partial_{y_i}\varphi(y')} \ge 0\\
    &G(\tilde{\delta}') \ge
    - (\tilde{\delta}' - \tilde{\delta})G'(\tilde{\delta}')
    + G(\tilde{\delta})
\end{align}
For two the geometric weighted mean trading function in equation
\ref{eq: weighted_geo_mean}, an analytic form of $F$ and $G$ has been derived in
equations \ref{eq: gm_fef}, \ref{eq: gm_ref},
\begin{align}
&F(\delta)= y_j\left(
    1 - \frac{y_i^{w_i/wj}}{\left(y_i+\gamma \delta\right)^{w_i/w_j} }
\right) \\
&G(\tilde{\delta}) = \frac{y_i}{\gamma}\left(
\frac{y_j^{w_j/w_i}}{\left(y_j - \tilde{\delta} \right)^{w_j/w_i}} - 1
\right),
\end{align}
for $\tilde{\delta} \le y_j$ and else $G(\tilde{\delta}) = \infty$ (invalid
trade).
For the weighted sum in equation \ref{eq: weighted_sum} the forward and
reverse exchange functions are
\begin{align}
    F(\delta) = \gamma \frac{p_i}{p_j}\delta \\
    F(\tilde{\delta}) = \frac{1}{\gamma} \frac{p_j}{p_i}\tilde{\delta},
\end{align}
for $\tilde{\delta} \le y_j$ and else $G(\tilde{\delta}) = \infty$ (invalid
trade).
In the case there is no analytical expression for the forward and reverse
exchange function, since the trade is uniquely defined by one variable either
$\delta$ or $\tilde{\delta}$, one can be fixed to compute a solution to the
trade acceptance condition numerically \ref{eq: trade_acceptance_condition}.
Recall the Newton method for a continuously differentiable function in
$\mathbb{R}$ with an initial guess $x_0$ near a root $a$ the iteration
\begin{align}
    x_{n+1} = x_n - \frac{f(x_n)}{f'(x_n)},
\end{align}
converges to an $a$. We can frame the trade acceptance condition as a
function of $\tilde{\delta}$ or $\delta$ the other respectively to arrive at
the iterations
\begin{align}
    \delta_{k+1} = \delta_k - \frac{\varphi(y+\gamma\delta_k e_i
    -\tilde{\delta}e_j) -\varphi(y)}{\partial_{y_i}\varphi(y+\gamma\delta_k e_i
-\tilde{\delta}e_j)}\label{eq: iter_backward},\\
\tilde{\delta}_{k+1} = \tilde{\delta}_k + \frac{\varphi(y+\gamma\delta e_i
    -\tilde{\delta}_ke_j) -\varphi(y)}{\partial_{y_j}\varphi(y+\gamma\delta e_i
-\tilde{\delta}_k e_j)}.
\end{align}
To solve for $\tilde{\delta}^{*}$, fix $\delta$ and we can choose the
iterates from the $\varepsilon$ ball defined by the Lipschitz continuity of
the derivative. Since $\varphi$ is continuously differentiable, the
derivative is also Lipschitz continuous, i.e. there exists a constant $L > 0$
such that
\begin{align}
   |f'(\tilde{\delta}) - f'(\tilde{\delta}^{*})|  \le L |\tilde{\delta} -
   \tilde{\delta}^{*} |.
\end{align}
Here $f$ is the function we are solving for $f(\tilde{\delta}) = \varphi(y
+\gamma \delta e_i -\tilde{\delta}e_j) - \varphi(y)$. Then
\begin{align}
    \varepsilon := \frac{\partial_{y_j}\varphi(y_j - \tilde{\delta}^{*}
    )}{2L}.
\end{align}
Assuming that $f'(\tilde{\delta}^{*})=\partial_{y_j}\varphi(y_j -
\tilde{\delta}^{*}) \neq 0$, we ensure that by concavity and monotonicity
that the $f'(\tilde{\delta})\neq 0$ for all $\delta$ in the ball of
$\varepsilon$, and also that
\begin{align}
    \frac{1}{\partial_{y_j}\varphi(y_j - \tilde{\delta}^{*})} \le
    \frac{1}{\partial_{y_j}\varphi(y_j - \tilde{\delta})} \le c,
\end{align}
for a $c\ge0$ where $|\tilde{\delta} - \tilde{\delta}^{*}| = \varepsilon$ . Thus
\begin{align}
    |f'(\tilde{\delta}) - f'(\tilde{\delta}^{*})| \le \frac{1}{2}
    |f'(\tilde{\delta}^{*})|
\end{align}
holds for all $\tilde{\delta}$ in the epsilon ball. Then by the mean value
theorem and Lipschitz continuity for all $\tilde{\delta}$ in the
$\varepsilon$ ball
\begin{align}
    &|f(\tilde{\delta}) - f(\tilde{\delta}^{*}) -
    f'(\tilde{\delta})(\tilde{\delta} - \tilde{\delta}^{*})| \le\\
    &\le |f(\tilde{\delta}) - f(\tilde{\delta}^{*}) -
    f'(\tilde{\delta})(\tilde{\delta} - \tilde{\delta}^{*})| +
    |f'(\tilde{\delta}) - f'(\tilde{\delta}^{*})|
    |\tilde{\delta}-\tilde{\delta}^{*} | \\
    &\le \int_{0}^{1} |f'(\tilde{\delta}^{*} + t(\tilde{\delta}
    -\tilde{\delta}) - f'(\tilde{\delta}^{*})| |\tilde{\delta} -
    \tilde{\delta}^{*} | dt
    +|f'(\tilde{\delta}) - f'(\tilde{\delta}^{*})|
    |\tilde{\delta}-\tilde{\delta}^{*} | \\
    &\le \int_0^1 L(1-t)
    |\tilde{\delta} - \tilde{\delta}^{*} |^{2}dt
    +|f'(\tilde{\delta}) - f'(\tilde{\delta}^{*})|
    |\tilde{\delta}-\tilde{\delta}^{*} | \\
    &\le \frac{1}{4c} |\tilde{\delta} - \tilde{\delta}^{*}| + \frac{1}{4c}
    |\tilde{\delta} - \tilde{\delta}^{*}| = \frac{1}{2c} |\tilde{\delta} -
    \tilde{\delta}^{*}|.
\end{align}
For a $\tilde{\delta}_0 $ in the epsilon ball the first iterate is also in
the ball
\begin{align}
    |\tilde{\delta}_{1} -\tilde{\delta}^{*} | &\le \left|\tilde{\delta}_0 -
    \frac{f(\tilde{\delta})}{f'(\tilde{\delta)}} - \tilde{\delta}^{*}
    \right| \\
    &\le |f'(\tilde{\delta}_0)|^{-1} |
    f(\tilde{\delta}_0) - f(\tilde{\delta}^{*}) -
    f'(\tilde{\delta}_0)(\tilde{\delta}_0 - \tilde{\delta}^{*})| \\
    &\le \frac{1}{2} |\tilde{\delta}_0-\tilde{\delta}^{*} |
\end{align}
Repeating this argument $k\ge0$ times we get
\begin{align}
    |\tilde{\delta}_{k+1}-\tilde{\delta}^{*} | \le \left( \frac{1}{2} \right)^{k}
    |\tilde{\delta}_0-\tilde{\delta}^{*}|
\end{align}
Thus the iterates of the Newton method $(\tilde{\delta}_k)_{k\in\mathbb{N}}$
lie in the $\varepsilon$ ball, the iteration is well defined since the
derivatives are non zero and the iterates converge to $\tilde{\delta}^{*}$ as
$k\to\infty$. To solve for $\delta$, fix $\tilde{\delta}$, the proof stays
the same.

\section{General \& Optimal Trades on one CFMM}
For general trades the agent choose a trade $(x, \tilde{x})$ from an
optimization problem that maximizes the change in the traders' portfolio $x
-\tilde{x}$ through a utility function. Where the agent prefers a change in
holding that maximizes a utility function $U: \mathbb{R}^{n}  \to \mathbb{R}
\cup \{-\infty\} $. For two trades $(x, \tilde{x}$ and $(x', \tilde{x}')$ the
agent prefers the trade $(x, \tilde{x})$ if
\begin{align}
    U(\tilde{x} -x) > U(\tilde{x}' - x').
\end{align}
And when $U \to -\infty$ the trade does not pass the trade acceptance
condition or the change in the portfolio is absolutely not preferred. The
utility function is chosen to be increasing and concave, and thereby differentiable.

\section{Utility optimization}
The trade that maximizes the utility can be framed as an optimization problem
\begin{align}
    \text{max} \quad &U(\tilde{x} - x) \label{eq: gen_trade}\\
    \text{s.t.:}\quad & \varphi(y + \gamma x - \tilde{x}) =
    \varphi(y)\nonumber \\
      &x\ge 0 \qquad \tilde{x} \ge 0\nonumber
\end{align}
It would be preferred if the problem in equation \ref{eq: gen_trade} is
convex, and can be solved efficiently fast using convex methods. Taking a
closer look at set
\begin{align}
    D := \left\{ x \in \mathbb{R}^{n}_+,\; \tilde{x} \in \mathbb{R}^{n}_+\;
    |\;
    \varphi(y + \gamma x - \tilde{x}) - \varphi(y) =0
\right\}
\end{align}
For $(x_1, \tilde{x}_1), (x_2, \tilde{x}_2) \in D$ and $\lambda \in [0, 1]$
reveals for $\left(\lambda x_1 + (1-\lambda)x_2, \lambda \tilde{x}_1
+(1-\lambda)\tilde{x}_2\right)$ that
\begin{align}
    &\varphi\Big( y+\gamma(\lambda x_1 + (1-\lambda)x_2) - (\lambda
    \tilde{x}_1 + (1-\lambda)x_2 \Big) - \varphi(y) =\\
    &=
    \varphi\Big(y + \lambda(\gamma x_1 - \tilde{x}_1) + (1-\lambda)(\gamma
    x_2 - \tilde{x}_2)\Big) - \varphi(y),\label{eq: convexity_trade_condition}
\end{align}
cannot be reduced further, thus the convex combination does belong to $D$
generally. But considering the relaxation from equality in $D$ to inequality
$\ge$, i.e.
\begin{align}
    D^{r} := \left\{ x \in \mathbb{R}^{n}_+,\; \tilde{x} \in \mathbb{R}^{n}_+\;
    |\;
    \varphi(y + \gamma x - \tilde{x}) - \varphi(y) \ge 0
\right\},
\end{align}
together with the concavity of $\varphi$ and $y = \lambda y + (1-\lambda)y$,
equation \ref{eq: convexity_trade_condition} can be reduced further
\begin{align}
    &\varphi\Big(y + \lambda(\gamma x_1 - \tilde{x}_1) + (1-\lambda)(\gamma
    x_2 - \tilde{x}_2)\Big) - \varphi(y) \\
    &=
    \varphi\Big(\lambda(y + \gamma x_1 - \tilde{x}_1) + (1-\lambda)(y - \gamma
    x_2 - \tilde{x}_2)\Big) - \varphi(y) \\
    &\ge \lambda \varphi(y + \gamma x_1 - \tilde{x}_1)
    + (1-\lambda)\varphi(y+\gamma x_2 - \tilde{x}_2) - \varphi(y).
\end{align}
For $(x_1, \tilde{x}_1), (x_2, \tilde{x}_2) \in D_{c}$ it is
\begin{align}
    \lambda \varphi(y + \gamma x_1 - \tilde{x}_1)
    + (1-\lambda)\varphi(y+\gamma x_2 - \tilde{x}_2) - \varphi(y)\ge 0,
\end{align}
hence $D^{r}$ is convex. So the relaxation of problem \ref{eq: gen_trade},
i.e.
\begin{align}
    \text{max} \quad &U(\tilde{x} - x) \label{eq: convex_gen_trade}\\
    \text{s.t.:}\quad & \varphi(y + \gamma x - \tilde{x}) \ge
    \varphi(y)\nonumber \\
      &x\ge 0 \qquad \tilde{x} \ge 0\nonumber
\end{align}
can be solved using convex methods. It is remains to show that any solution
of the convex relaxation \ref{eq: convex_gen_trade} is a feasible solution to
the original \ref{eq: gen_trade}. Trivially if a solution of the relaxation
\ref{eq: convex_gen_trade} satisfies the equality condition $\varphi(y +
\gamma x - \tilde{x}) = \varphi(y)$ then it is also solution to the original
\ref{eq: gen_trade}. In the case of strict inequality
\begin{align}
    \varphi(y + \gamma x - \tilde{x}) > \varphi(y).
\end{align}
Since $\varphi$ is increasing in $x$ and decreasing in $\tilde{x}$, we
can decrease $x$ by $\varepsilon>0$ and increase $\tilde{x}$ by
$\tilde{\varepsilon}>0$ until
\begin{align}
    \varphi(y + \gamma (x-\varepsilon) - (\tilde{x}+\tilde{\varepsilon})) \ge \varphi(y).
\end{align}
But then the utility of $(x-\varepsilon, x+\tilde{\varepsilon})$ is bigger
than that of $(x,\tilde{x})$. Hence, $(x,\tilde{x})$ is not an optimal
solution.

\subsection{Optimality Condition\label{sec: optimality condition}}
The KKT-conditions of the minimization of
equation \ref{eq: convex_gen_trade}, i.e.
\begin{align}
    \text{min} \quad &-U(\tilde{x}-x)\\
    \text{s.t.:}\quad & (x,\tilde{x}) \in D^{r},
\end{align}
exists, because the convex optimization problem satisfies the Slater CQ since
$\varphi$ increasing we can set $\tilde{x}=$ and any $x>0$ such that. \begin{align}
    \varphi(y) - \varphi(y + \gamma x) < 0.
\end{align}
Since Slater CQ is satisfied so is Abadie CQ, then there is a KKT point
$(x^{*} , \tilde{x}^{*}, \lambda^{*} , \omega^{*} , \mu^{*} )$, which is a
sufficient condition for optimality in the convex case. So $(x^{*} ,
\tilde{x}^{*}$) is a local minimum of the optimization problem.
Thereby any $(x, \tilde{x}, \lambda, \omega, \mu)$ satisfying
\begin{align}
    \begin{cases}
    \nabla_x L(x, \tilde{x}, \lambda, \omega, \mu) = 0,\\
    \nabla_{\tilde{x}} L(x, \tilde{x}, \lambda, \omega, \mu) = 0,\\
    \lambda \ge 0,\; g(x, \tilde{x}) \le 0,\; \lambda g(x, \tilde{x}) = 0,\\
    \omega \ge 0,\; \omega x = 0,\\
    \mu \ge 0,\; \mu \tilde{x} = 0\\
    \end{cases},
\end{align}
is a KKT point.
Here $L:
\mathbb{R}^{n}\times\mathbb{R}_+\times\mathbb{R}_+^{n}\times\mathbb{R}_+^{n}
\to \mathbb{R} $ is the Lagrangian, and $\lambda \in \mathbb{R}_+, \omega \in
\mathbb{R}_+^{n}, \mu \in \mathbb{R}_+^{n}$  are the Lagrange multipliers and
$g(x, \tilde{x}) = \varphi(y) - \varphi(y + \gamma x - \tilde{x})$ just a
notational convention. Then the Lagrangian is
\begin{align}
    L(x, \tilde{x},\lambda, \omega, \mu) = -U(\tilde{x}-x) + \lambda g(x,
    \tilde{x}) - \omega x - \mu \tilde{x},
\end{align}
and the first two conditions give
\begin{align}
    0 = \nabla_x L = \nabla U(\tilde{x}-x) - \lambda \gamma \nabla\varphi(y') -
    \underbrace{\omega}_{\ge0},\\
    0 = \nabla_{\tilde{x}} L = - \nabla U(\tilde{x}-x) + \lambda  \nabla\varphi(y') -
    \underbrace{\mu}_{\ge0},
\end{align}
where $y' = y-\gamma x + \tilde{x}$. The two inequalities give
\begin{align}
    \lambda \gamma \nabla \varphi(y') \le \nabla U(\tilde{x}-x)\le\lambda
    \nabla\varphi(y').
\end{align}
Specific interpretation for this condition is when the trade is not
increasing the utility, e.g. for a solution $x = \tilde{x} = 0$. Then the
condition, after dividing with $\lambda P_n$ becomes
\begin{align}
    \gamma p \le \frac{1}{\lambda P_n}\nabla U(0) \le p.
\end{align}
The set of prices $p$ which satisfy this condition can is denoted by
\begin{align}
    K := \left\{p \in \mathbb{R}_+^{n}| \;\exists \alpha\ge0:\; \gamma p \le \alpha U(0) \le p \right\}
\end{align}
is in literature referred to as the no-trade cone for the utility U for trade
function $\varphi$. The set is a cone since multiplying by any $p \in K$ by
an $\eta > 0$, by division $\frac{\alpha}{\eta}$ is still bigger or equal to
zero. The Set $K$ is also convex, for any $p_1, p_2 \in K$ and $\eta \in [0,
1]$, $\eta p_1 + (1-\eta)p_2 \in K$, because the condition is additive and a
$K$ a cone. The cone is nonempty since $0 \in K$ with the choice of $\alpha =
0$.

The case of linear utility $U(\psi) = z^{T} \psi$, with $z \in
\mathbb{R}_+^{n}$, then
\begin{align}
    \gamma p \le \frac{1}{\lambda P_n}z \le p,
\end{align}
component wise naturally.
\section{Multi-DEX Trades}
The agent might decide to trade with multiple DEXs, hence the trade will
interact with multiple CFMMs. In light of multi asset trades
\ref{eq: convex_gen_trade} a generalization to an interaction with multiple CFMMS can
be made. Consider assets $j \in \{1, \ldots,n\}$, labeled as such. Consider
CFMMs $i \in \{1,\ldots,m\}$, each trading $N_i \subseteq \{1,\ldots,n\}  $
assets which are specifically ordered based on the construction of the CFMM
$i$ and trading function $\varphi_i: \mathbb{R}^{|N_i|}_+\to \mathbb{R}$ with
liquidity $y_i$ and trading fee $\gamma_i$. To link the ordering of $N_i$ and
asset labels $i$, consider binary matrices $A_i \in \mathbb{R}^{n\times
|N_i|}$ with
\begin{align}
    (A_i)_{kj} =
    \begin{cases}
        1 & \text{for } k \in N_i, \; k = j \\
        0 & \text{else}.
        \end{cases}
\end{align}
The agent trades with each CFMM $i$ by a trade $(x_i, \tilde{x}_i)$, the
mapping to the asset indices $j$ is done through the matrices $A_i$. Hence,
the global trade vector interacting with all CFMMs is $\psi \in
\mathbb{R}^{n}$ is defined as
\begin{align}
    \psi = \sum_{i=1}^{m} A_i(\tilde{x}_i-x_i).
\end{align}
The implication if $\psi_k > 0$, means that the agent does not give any
amount of asset $k$ to the set of CFMMs. From this the agent still might
trade asset $k$, but receives as much as given. Analogous to multi-asset
trades, the agent judges a trade with a utility function
\begin{align}
    U:\mathbb{R}^{n}&\to\mathbb{R}\\
    \psi &\mapsto U(\psi),
\end{align}
where $U$ is concave, differentiable and increasing. In the case of $U(\psi)
= -\infty$ the trade is not acceptable, and the agent does not trade. E.g. if
the agent wants to increase its holdings $h\in \mathbb{R}_+^{n}$ after the
trade, then $U(\psi) = \psi + h \ge 0$ is a constraint satisfying the
condition of increasing the holdings. For any trade which does not satisfy
the constraint, the utility is $-\infty$, hence not acceptable.

A set of optimal trades with CFMMs maximizing the utility is
\begin{align}
    \text{max} \quad &U\left(\sum_{i=1}^{m} A_i(\tilde{x}_i-x_i)\right) \label{eq: multi_dex_trade_orig}\\
    \text{s.t.:}\quad
                      & (x_i, \tilde{x}_i) \in D_i \;\; \forall i \in
                      \{1,\ldots,m\},
\end{align}
where
\begin{align}
D_i = \left\{(x, \tilde{x})|\; x, \tilde{x} \in \mathbb{R}^{|N_i|}_+, \;\;
    \varphi_i(y_i+\gamma_ix - \tilde{x}) = \varphi_i(y_i)
\right\}
\end{align}
It is already shown that the convex relaxation of the optimization problem
satisfies the solution to the original, the proof stays the same
component wise for all CFMMs $i \in \{1,\ldots,m\}$. The convex relaxation
problem is
\begin{align}
    \text{max} \quad &U\left(\sum_{i=1}^{m} A_i(\tilde{x}_i-x_i)\right) \label{eq: multi_dex_trade}\\
    \text{s.t.:}\quad & (x_i, \tilde{x}_i) \in D_i^{r}
          \;\; \forall i \in
          \{1,\ldots,m\},
\end{align}
where
\begin{align}
D_i^{r}  = \left\{(x, \tilde{x})|\; x, \tilde{x} \in \mathbb{R}^{|N_i|}_+, \;\;
    \varphi_i(y_i+\gamma_ix - \tilde{x}) \ge \varphi_i(y_i)
\right\}
\end{align}
Similar to the argumentation in section \ref{sec: optimality condition} the
Slater CQ is satisfied for all CFMM $\varphi_i$ $i \in \{1,\ldots,m\}$ by
setting the outgoing amount to zero. This is a sufficient condition for
optimality in the convex case. The Lagrangian of the problem is
\begin{align}
    L(x, \tilde{x}, \nu, \lambda, \omega, \mu) = -U(x, \tilde{x}) +
    \nu^{T}\left(  \psi - A(x-\tilde{x})\right) + \lambda^{T} g(x, \tilde{x})
    -\omega^{T} x - \mu^{T} \tilde{x}.
\end{align}
The KKT conditions give
\begin{align}
    &0 = \nabla L = -\nabla U(x, \tilde{x}) + \nu = 0 \\
    \implies& \nabla U(x, \tilde{x}) = \nu.
\end{align}
Also for all $i \in \{1,\ldots,m\}$
\begin{align}
    &\nabla_{x_i} L = A_i^{T}\nabla U(x, \tilde{x}) - \lambda_i\gamma\nabla\varphi_i(y_i
    +\gamma_ix_i -\tilde{x}_i) - \underbrace{\omega}_{\ge0}\\
    &\nabla_{\tilde{x}_i} L = -A_i^{T}\nabla U(x, \tilde{x}) + \lambda_i\nabla\varphi_i(y_i
    +\gamma_ix_i -\tilde{x}_i) - \underbrace{\omega}_{\ge0}.
\end{align}
Thus,
\begin{align}
    &A_i^{T}\nu \ge \lambda_i\gamma\nabla\varphi_i(y_i
    +\gamma_ix_i -\tilde{x}_i)\\
    &A_i^{T}\nu \le \lambda_i\nabla\varphi_i(y_i
    +\gamma_ix_i -\tilde{x}_i).
\end{align}
Combining gives
\begin{align}
    \lambda_i\gamma\nabla\varphi_i(y_i
    +\gamma_ix_i -\tilde{x}_i)\le
    A_i^{T}\nu \le \lambda_i \nabla\varphi_i(y_i
    +\gamma_ix_i -\tilde{x}_i).
\end{align}
When the optimal trade is zero, the agent does not trade. This is the
case when $\psi =0$ and $x_i=\tilde{x}_i=0$ for all $i \in \{1,\ldots,m\}$.
Thus, the no-trade cones are defined through the following inequality
component-wise
\begin{align}
    \gamma_i\lambda_i p_i \le A_i^{T} \nabla U(0)\le\lambda p_i \qquad
    \forall i \in \{1,\ldots,m\},
\end{align}
where $p_i = \nabla \varphi_i(y_i)$ are the marginal exchange rates of CFMM
$i$ at liquidity $y_i$.
\subsection{Transaction Cost}
The mechanics of a blockchain transaction include the agent paying for
transaction fees, depending on the computational cost of the transaction.
E.g. for any trade on the Ethereum Virtual Machine (EVM) \cite{evm_optcodes},
the transaction cost (gas fee) is calculated based on assembly opt-codes
called during the transaction. In light of this let $q_i \in \mathbb{R}_+$
some cost related to executing a trade with CFMM $i$. The cost $q_i$ is only
given if the transaction is executed (or partially executed), i.e. $\tilde{x}_i -
x_i \neq 0$. To keep track with which CFMM the agent trades, define a binary
vector $\eta \in \{0, 1\}^{m} $ where
\begin{align}
    \eta =
    \begin{cases}
        &1\quad \text{nonzero trade with CFMM $i$} \\
        &0\quad \text{else}.
    \end{cases}
\end{align}
The total cost of conducting a general trade with $m$ CFMM is $q^{T}\eta$.
Assuming there is a maximum $x_i^{*} \in \mathbb{R}_+^{|N_i|}$. It is
necessary that the trades on CFMM $i$ do not give more amount of assets that
the maximum in a nonzero trade, i.e. $x_i \le \eta_ix_i^{*}$. The optimal
trade needs to maximize the utility of the trade minus the transaction cost
\begin{align}
    \text{max} \quad &U\left(\sum_{i=1}^{m} A_i(\tilde{x}_i-x_i)\right) - q^{T}\eta  \label{eq: trade_transaction_cost}\\
    \text{s.t.:}\quad & (x_i, \tilde{x}_i) \in D_i^{r}\\
          &x_i \le \eta_i x_i^{*}, \quad \forall i \in\{1,\ldots,m\} \\
          &\eta_i \in \{0, 1\},
\end{align}
with variables $x_i, \tilde{x}_i$, $i\in\{1,\ldots,m\}$ and $\eta \in
\{0,1\}^{m} $. This is a mixed integer convex optimization problem. A
possible solution might evolve the relaxation to the condition $\eta \in [0,
1]^{m}$ to transform it to a convex optimization problem. And introduce a
learned parameter $\tau(\theta)$ with some weights $\theta$, as a threshold
to flip the value of $\eta$ to one or zero. This however would evolve
learning $\tau$, by solving the mixed integer problem many times in real
world scenarios. The trade-off would be to have a faster solution when
computing in real time.
\subsection{Cyclic Arbitrage}
Going back to the general trading problem without transaction fees in
equation \ref{eq: multi_dex_trade}. It is important to note that the solution
does specify in which order the trades should be executed and that the
computational effort will grow larger with more CFMM. Hence, in a large space
of CFMM it is a separate problem of finding cyclic arbitrage path. This path
will give the order of the trades and specify for which CFMM the optimization
problem needs to be framed, dropping the rest.

Consider now that the agent has found a cyclic arbitrage path of $m \in
\mathbb{N} $ CFMM. For an arbitrage trade the agent needs the trade to be
valid, i.e. $\psi \neq 0$. The agent also needs to have $\psi \ge 0$, since
at the end of a set of trades, the arbitrage trade receives a positive amount
of at least one asset. This hold only for $\psi$, i.e. at the end of the
trade. The utility function needs to be bound in the domain
$\mathbb{R}_+^{n}$, and there is an arbitrage if there is a nonzero solution.
An optimal arbitrage is one that maximizes the utility under the constraints
i.e. $U(\psi^{*}) > U(0)$ with $\psi^{*}$ the optimal solution. The
optimization problem is then
\begin{align}
    \text{max} \quad &U\left(\sum_{i=1}^{m} A_i(\tilde{x}_i-x_i)\right) \label{eq: arb_trade}\\
    \text{s.t.:}\quad & (x_i, \tilde{x}_i) \in D_i^{r}
          \;\; \forall i \in
          \{1,\ldots,m\},
\end{align}
Again the solution where $\sum_{i=1}^{m} A_i(\tilde{x}_i-x_i)=0$ is valid, but in this case there is no arbitrage
and the agent does not trade.
\subsubsection{Example}
A common case is an $n$ cyclic arbitrage under linear utility
\begin{align}
 U(x, \tilde{x}) =1^{T}\left(\sum_{i=1}^{m} A_i(\tilde{x}_i-x_i)\right) ,
\end{align}
with $n$ CFMMs $\varphi_i(y_i) =
\sqrt{y_{i,1}y_{i,2}}$, with $n$ distinct assets via
\begin{align}
    1 \xrightarrow{\varphi_1} 2 \xrightarrow{\varphi_2} 3
    \xrightarrow{\varphi_3} \hdots \xrightarrow{\varphi_{n-1}} n
    \xrightarrow{\varphi_n} 1.
\end{align}
W.l.o.g. for each trade let $(x_i, \tilde{x}_i) = (\delta_1e_1,
\tilde{\delta}e_2)$, where $e_1, e_2$ are the basis vectors in
$\mathbb{R}^{2}$. Then the matrices $A_i \in \mathbb{R}^{n\times2}  $ take the shape
\begin{align}
   A_1 =
   \begin{pmatrix}
   1 & 0 \\
   0 & 1 \\
   \vdots & \vdots \\
   0 & 0
    \end{pmatrix},\;
   A_2 =
   \begin{pmatrix}
   0 & 0 \\
   1 & 0 \\
   0 & 1 \\
   \vdots & \vdots \\
   0 & 0
    \end{pmatrix},\; \ldots  \;
   A_{n-1} =
   \begin{pmatrix}
   0 & 0 \\
   \vdots & \vdots \\
   1 & 0\\
   0 & 1
    \end{pmatrix},\;
   A_n =
   \begin{pmatrix}
   0 & 1 \\
   \vdots & \vdots \\
   1 & 0
    \end{pmatrix}.
\end{align}
Also, the general trade vector is
\begin{align}
    \sum_{i=1}^{m} A_i(\tilde{x}_i-x_i) =
    \begin{pmatrix}
        \tilde{\delta}_n - \delta_1\\
        \tilde{\delta}_1 -\delta_2\\
        \vdots\\
        \tilde{\delta}_{n-1} - \delta_n
    \end{pmatrix}.
\end{align}
The convex constraints component wise can be written as
\begin{align}
    y_{i,1}y_{i,2} - (y_{i,1}+\gamma_i \delta_i)(y_{i,2}-\tilde{\delta}_i)
    \le 0.
\end{align}
This can be rewritten into a quadratic form by simplifying
\begin{align}
    &y_{i,1}y_{i,2} - (y_{i,1}+\gamma_i \delta_i)(y_{i,2}-\tilde{\delta}_i) \\
    &=y_{i,1}\tilde{\delta}_i - \gamma_i y_{i,2}\delta_i +
    \gamma_i\delta_i\tilde{\delta}_i \\
    & =
    \begin{pmatrix} -\gamma_i y_{i,2} & y_{i,1} \end{pmatrix}
    \begin{pmatrix} \delta_i \\ \tilde{\delta}_i \end{pmatrix}
    +\gamma_i
    \begin{pmatrix} \delta_i & \tilde{\delta}_i \end{pmatrix}
   \begin{pmatrix} 0 & 1 \\ 0 & 0 \end{pmatrix}
    \begin{pmatrix} \delta_i \\ \tilde{\delta}_i \end{pmatrix}.
\end{align}
Set $\xi_i = \begin{pmatrix} \delta_i & \tilde{\delta}_i\end{pmatrix} ^{T}
\in \mathbb{R}^{2}$, $c_i = \begin{pmatrix} -\gamma_iy_{i,2}&
    y_{i,1}\end{pmatrix}^{T} \in \mathbb{R}^{2}$ and $B = \begin{pmatrix} 0 &
1\\ 0 & 0\end{pmatrix}$. Then the constraint can be written quadratic
constraint
\begin{align}
    c\xi_i + \left(\sum_{i=j}^{m} A_j(\tilde{x}_j-x_j)\right)_i^{T} A\xi_i \le 0.
\end{align}
Together with $R = \begin{pmatrix} -1 & 0 \\ 0 & 1 \end{pmatrix}$ the
optimization problem becomes a quadratically constrained linear problem.
\begin{align}
    \text{max} \quad & \sum_{i=1}^{n}1^{T}A_iR\xi_i \label{eq: quadratically_constrained}\\
    \text{s.t.:}\quad
          & \sum_{i=1}^{n}A_iR\xi_i \ge 0 \\
          & \xi_i \ge 0, \;\;
          c_i^{T} \xi_i + \gamma_i \xi_i^{T} B\xi_i \le 0
          \;\; \forall i \in \{1,\ldots,n\}.
\end{align}
Additionally, it is possible use matrix factorization to factor $\sum_{i=1}^{n}A_iR\xi_i$
into $A\xi$ for $A \in \mathbb{R}^{n\times2n}$ with entries $-1, 0,
1$, with $\xi = \begin{pmatrix} \delta_1 & \tilde{\delta}_1 & \hdots &
\delta_n & \tilde{\delta}_n \end{pmatrix}^{T}  \in \mathbb{R}^{2n}  $. Also,
it is possible to pull the inequality constraint into a single constraints
under the same $\xi$ with a $c \in \mathbb{R}^{2n}$ and $D \in \mathbb{R}^{2n
\times 2n}$ consisting of $n$ blocks of $\gamma_iB$, to get
\begin{align}
    \text{max} \quad & 1^{T}A\xi \label{eq: quadratically_constrained_general}\\
    \text{s.t.:}\quad & c\odot \xi + \xi \odot D\xi \le 0\\
          & \xi \ge 0,
\end{align}
where $\odot$ denotes the componentwise multiplication
